\section*{Capítulo \ref{ch:b2:pt-paths} --- Termobarometria e trajetórias P--T--t}

\begin{solution}{exo:b2:pt-paths:1}
(a) A zona da clorita, B zona da biotita, C zona da granada, D zona da estaurolita, E zona da cianita,
F zona da silimanita. (b) Cinco: as isógradas da biotita, da granada, da estaurolita, da cianita e da silimanita,
cada uma entre duas amostras sucessivas, com o grau subindo de A para F. (c) Uma isógrada
diz apenas que uma reação ocorreu; um gradiente de campo precisa de pressões e temperaturas, que
vêm de reações calibradas ou da termobarometria da assembleia de pico de cada
rocha.
\end{solution}

\begin{solution}{exo:b2:pt-paths:2}
(a) $(X_{\text{Mg}}/X_{\text{Fe}})_{\text{grt}} = 0.10/0.80 = 0.125$,
$(X_{\text{Mg}}/X_{\text{Fe}})_{\text{bt}} = 0.45/0.50 = 0.900$; $K_D = 0.1389$, $\ln K_D =
-1.9741$. (b) $3R = 24.942$; denominador $20.19 + 24.942\times1.9741 = 69.428$; a
\qty{5}{kbar}, $T = 54185/69.428 = \qty{780.5}{K} = \qty{507}{\celsius}$; a
\qty{7}{kbar}, $T = (52630 + 2177)/69.428 = \qty{789.4}{K} = \qty{516}{\celsius}$.
(c) Dois quilobares a mudam em apenas \qty{9}{K}: um termômetro.
\end{solution}

\begin{solution}{exo:b2:pt-paths:3}
(a) $z = P/(\bar\rho g) = 6.0\times10^8/(2800\times9.81) = \qty{21.8}{km}$.
(b) $(580 - 10)/21.8 = \qty{26}{\celsius\per\km}$. (c) A \qty{580}{\celsius}, a linha
fica em $3.67 + 0.8\times2.10 = \qty{5.35}{kbar}$; a \qty{6.0}{kbar}, a rocha está do
lado de alta pressão: cianita.
\end{solution}

\begin{solution}{exo:b2:pt-paths:4}
(a) Reagentes $2\times140.693 + 3\times18.015 = \qty{335.431}{g}$; produtos $277.112 +
58.320 = \qty{335.432}{g}$: equilibrada dentro do arredondamento. (b) $2\times43.65 =
\qty{87.30}{cm^3}$ se tornam $107.50 + 24.63 = \qty{132.13}{cm^3}$: $+\qty{51.4}{\percent}$.
(c) Olivina $281.386/87.30 = \qty{3.223}{g.cm^{-3}}$; produtos $335.432/132.13 =
\qty{2.539}{g.cm^{-3}}$. (d) A rocha incha metade de seu volume e fica mais leve; ela
também armazena água em seus minerais.
\end{solution}

\begin{solution}{exo:b2:pt-paths:5}
(a) Médias $-1.206$ a \qty{1072.15}{K} e $-1.782$ a \qty{823.15}{K}. Com $1/T =
9.3271\times10^{-4}$ e $1.21485\times10^{-3}$, $a = (-1.206 + 1.782)/(9.3271\times10^{-4}
- 1.21485\times10^{-3}) = -2041.5$ K e $b = -1.206 + 2041.5/1072.15 = 0.6982$.
(b) $\Delta_r H + 2070\times0.311 = 2041.5\times24.942 = \qty{50920}{J}$, logo $\Delta_r H =
\qty{50276}{J}$; $\Delta_r S = 0.6982\times24.942 = \qty{17.41}{J.K^{-1}}$. (c) A
\qty{5}{kbar}: para $K_D = 0.18$, 588.1 contra \qty{587.5}{\celsius}; para $0.10$, 419
contra \qty{425}{\celsius}; para $0.30$, 819 contra \qty{806}{\celsius}. Dentro do intervalo
dos experimentos, as duas calibrações concordam, embora suas constantes difiram de
\qty{2.4}{kJ} e \qty{2.8}{J.K^{-1}}; os erros de $\Delta_r H$ e $\Delta_r S$
se compensam ali e crescem fora dele. Uma calibração não deve ser usada longe
das condições de seus experimentos.
\end{solution}

\begin{solution}{exo:b2:pt-paths:6}
(a) Troca: $(113.12 + 154.30) - (115.32 + 149.65) = \qty{+2.45}{cm^3} =
\qty{0.245}{J.bar^{-1}}$, contra 0.311. GASP: $125.28 + 2\times44.15 + 22.69 -
3\times100.79 = \qty{-66.10}{cm^3} = \qty{-6.61}{J.bar^{-1}}$, contra $-6.21$, o valor
da calibração a \qty{800}{\celsius} e \qty{10}{kbar}, onde a dilatação térmica e a
compressão mudam os volumes. Nos dois casos, a variação da GASP é de vinte a trinta vezes
a da troca. (b) $\partial T/\partial P = 0.311/62.96 = \qty{4.9e-3}{K/bar}$:
\qty{4.9}{K} por quilobar. (c) A inclinação da linha é $22.80 + 8.314\times3\ln 0.20/6.21
= 22.80 - 6.46 = \qty{16.3}{bar/K}$; \qty{17}{K} deslocam a pressão de
\qty{280}{bar}.
\end{solution}

\begin{solution}{exo:b2:pt-paths:7}
(a) $P_0 = 22.80\times600 - 1093 = \qty{12587}{bar}$; $3\ln(0.08/0.30) = -3.9653$; $RT =
8.314\times873.15 = \qty{7259.4}{J}$; $P = 12587 + 7259.4\times(-3.9653)/6.21 = 12587 -
4635 = \qty{7952}{bar}$. (b) $3\ln(0.08/0.35) = -4.4277$, $P = 12587 - 5176 =
\qty{7411}{bar}$: \qty{0.54}{kbar} a menos, de modo que o plagioclásio deve ser analisado com cuidado.
(c) $7952\times10^5/(2800\times9.81) = \qty{29.0}{km}$.
\end{solution}

\begin{solution}{exo:b2:pt-paths:8}
(a) $z = P/(\bar\rho g)$: \qty{118.9}{km}, \qty{34.0}{km} e \qty{15.3}{km}.
(b) $(118.9 - 34.0)/2.2 = \qty{38.6}{km/Myr} = \qty{3.9}{cm.yr^{-1}}$ e $(34.0 -
15.3)/1.1 = \qty{17.0}{km/Myr} = \qty{1.7}{cm.yr^{-1}}$, um pouco acima dos
3.4 e 1.6 publicados: sua conversão de pressão em profundidade dá diferenças de profundidade
um pouco menores. (c) $\sigma = \sqrt{0.9^2 + 0.9^2} = \qty{1.3}{Myr}$, logo $2.2 \pm 1.3$
milhões de anos e uma taxa entre $84.9/3.47 = 24$ e $84.9/0.93 = \qty{91}{km/Myr}$,
2.4 a \qty{9.1}{cm.yr^{-1}}: rápida em qualquer caso. (d) Mesmo a taxa mais baixa é cerca de cinquenta vezes
os \qty{0.5}{mm.yr^{-1}} do modelo: a erosão sozinha não pode fazê-lo.
\end{solution}

\begin{solution}{exo:b2:pt-paths:9}
(a) $\partial T/\partial\ln K_D = 3RT/(\Delta_r S - 3R\ln K_D) = 24.942\times860.6/62.96 =
\qty{341}{K}$; $\partial T/\partial P = 0.311/62.96 = \qty{4.9}{K}$ por quilobar.
(b) $\sigma_T = \sqrt{(341\times0.05)^2 + 4.9^2} = \sqrt{17.0^2 + 4.9^2} =
\qty{17.7}{K}$. (c) O erro analítico sobre $K_D$ domina; sozinho, ele vale dois terços dos
\qty{25}{\celsius} das melhores calibrações, que incluem também os erros da
própria calibração.
\end{solution}

\begin{solution}{exo:b2:pt-paths:10}
(a) Granada $0.10/0.75 = 0.1333$, biotita $0.42/0.52 = 0.8077$, $K_D = 0.1650$, $\ln K_D =
-1.8018$, denominador $20.19 + 24.942\times1.8018 = 65.130$; $3\ln(0.07/0.25) = -3.8189$.
A \qty{6000}{bar}: $T = 54496/65.130 = \qty{836.7}{K} = \qty{563.6}{\celsius}$, depois $P
= 22.80\times563.6 - 1093 + 8.314\times836.7\times(-3.8189)/6.21 = \qty{7479}{bar}$.
A \qty{7479}{bar}: $T = \qty{570.6}{\celsius}$, $P = \qty{7603}{bar}$. A
\qty{7603}{bar}: $T = \qty{571.2}{\celsius}$, $P = \qty{7614}{bar}$; mais uma vez, $T =
\qty{571.3}{\celsius}$. A rocha se equilibrou a cerca de \qty{571}{\celsius} e
\qty{7.6}{kbar}. (b) A \qty{571}{\celsius}, a linha cianita $=$ silimanita fica em $3.67
+ 0.71\times2.10 = \qty{5.16}{kbar}$: cianita. (c) O barômetro foi calibrado para a
cianita; com silimanita, a reação, e sua linha, são outras, já que os dois polimorfos não
têm a mesma energia de Gibbs.
\end{solution}

\begin{solution}{exo:b2:pt-paths:11}
(a) $\tau = L^2/(\pi^2\kappa)$: para \qty{35}{km}, $1.225\times10^9/(9.870\times10^{-6}) =
\qty{1.24e14}{s} = \qty{3.9}{Myr}$; para \qty{70}{km}, quatro vezes mais,
\qty{15.7}{Myr}. (b) Setenta milhões de anos são de quatro a dezoito vezes esses tempos: o calor
se propaga pela crosta espessada muito mais depressa do que a erosão a remove. (c) A rocha
atinge sua maior profundidade de uma vez, no empilhamento; o calor chega ao longo das
dezenas de milhões de anos seguintes, enquanto a erosão a traz lentamente para cima: a temperatura continua
subindo por algum tempo durante a descompressão. (d) Uma erosão em sete milhões de anos
traria a rocha para cima antes que a condução a tivesse aquecido: um pico térmico mais baixo, atingido
cedo, a pressão mais alta, e uma trajetória fria de volta à superfície.
\end{solution}

\begin{solution}{exo:b2:pt-paths:12}
(a) A crosta pesa $35000\times2800\times9.81 = \qty{0.961}{GPa}$; os
\qty{1.839}{GPa} restantes exigem $1.839\times10^9/(3300\times9.81) = \qty{56.8}{km}$ de manto: profundidade de
\qty{91.8}{km}, como os \qty{90}{km} ou mais do estudo original. (b) $(4.0 -
0.961)\times10^9/(3300\times9.81) = \qty{93.9}{km}$, profundidade de \qty{128.9}{km}.
(c) $700/91.8 = \qty{7.6}{\celsius\per\km}$: uma razão muito alta entre pressão e
temperatura, a série das zonas de subducção. (d) $(91.8 - 35)/34 = \qty{1.7}{Myr}$.
\end{solution}

\begin{solution}{pb:b2:pt-paths:1}
\noindent\textbf{Parte I.}

\textbf{1.} Núcleo da granada $0.080/0.760 = 0.1053$, biotita $0.40/0.58 = 0.6897$; $K_D =
0.1527$, $\ln K_D = -1.8793$.

\textbf{2.} Borda $0.120/0.720 = 0.1667$, biotita da matriz $0.46/0.53 = 0.8679$; $K_D =
0.1921$, $\ln K_D = -1.6497$.

\textbf{3.} $3R = 24.942$; numerador $52630 + 0.311\times6000 = 54496$; denominador $20.19
+ 24.942\times1.8793 = 67.064$; $T = \qty{812.6}{K} = \qty{539.5}{\celsius}$.

\textbf{4.} Denominador $20.19 + 24.942\times1.6497 = 61.337$; $T = 54496/61.337 =
\qty{888.5}{K} = \qty{615.3}{\celsius}$.

\textbf{5.} $0.311\times1000/67.064 = \qty{4.6}{K}$ para o núcleo e $0.311\times1000/61.337
= \qty{5.1}{K}$ para a borda.

\textbf{6.} Um termômetro se aplica a minerais que estiveram em equilíbrio juntos. A
biotita incluída no núcleo ficou isolada do resto da rocha quando o núcleo cresceu;
a biotita da matriz trocou por último com a borda. Emparelhar o núcleo com a matriz misturaria
dois momentos.

\noindent\textbf{Parte II.}

\textbf{7.} Núcleo $22.80\times539.5 - 1093 = \qty{11208}{bar}$; borda $22.80\times615.3 -
1093 = \qty{12936}{bar}$.

\textbf{8.} $3\ln(0.050/0.25) = -4.8283$; $RT = 8.314\times812.65 = \qty{6756.4}{J}$; $P =
11208 - 6756.4\times4.8283/6.21 = 11208 - 5253 = \qty{5955}{bar}$.

\textbf{9.} $3\ln(0.060/0.28) = -4.6213$; $RT = 8.314\times888.45 = \qty{7386.6}{J}$; $P
= 12936 - 5497 = \qty{7439}{bar}$.

\textbf{10.} A \qty{7439}{bar}: $T = (52630 + 2314)/61.337 = \qty{895.8}{K} =
\qty{622.6}{\celsius}$; a essa temperatura, $P = 13102 - 5542 = \qty{7560}{bar}$; depois $T
= \qty{623.2}{\celsius}$. A borda cresceu a cerca de \qty{623}{\celsius} e \qty{7.56}{kbar}.

\textbf{11.} A \qty{5954}{bar}, $T = \qty{539.2}{\celsius}$ e $P = \qty{5950}{bar}$: o
núcleo cresceu a cerca de \qty{539}{\celsius} e \qty{5.95}{kbar}. A primeira estimativa,
\qty{6}{kbar}, já estava próxima, e o termômetro se desloca apenas \qty{4.6}{K} por quilobar.

\textbf{12.} Núcleo $5950\times10^5/(2800\times9.81) = \qty{21.7}{km}$; borda
$7560\times10^5/27468 = \qty{27.5}{km}$.

\textbf{13.} A \qty{539}{\celsius}, a linha fica em $3.67 + 0.39\times2.10 =
\qty{4.49}{kbar}$, abaixo de \qty{5.95}{kbar}: cianita, coerente com a cianita incluída
no núcleo. A \qty{623}{\celsius}, ela fica em $5.77 + 0.23\times2.09 = \qty{6.25}{kbar}$,
abaixo de \qty{7.56}{kbar}: cianita de novo.

\noindent\textbf{Parte III.}

\textbf{14.} $+\qty{84}{\celsius}$ e $+\qty{1.6}{kbar}$ (cerca de \qty{6}{km} mais fundo): a
rocha estava sendo soterrada e aquecida, em sua \omterm{def:b2:pt-paths:paths}{trajetória progressiva}.

\textbf{15.} A \qty{623}{\celsius} ou mais, abaixo de \qty{6.25}{kbar}: depois que a borda cresceu, a
rocha perdeu pelo menos \qty{1.3}{kbar}, cerca de \qty{5}{km} de cobertura, permanecendo pelo
menos tão quente. Ela foi exumada ainda quente.

\textbf{16.} A clorita é um mineral hidratado de baixo grau: ela registra a \omterm{def:b2:pt-paths:paths}{trajetória
retrógrada} tardia, o resfriamento a pequena profundidade, com água.

\textbf{17.} Soterramento e aquecimento até a borda, descompressão sem resfriamento, depois resfriamento:
desenhada com a temperatura para a direita e a pressão para cima, a trajetória gira no sentido horário. O
pico de pressão veio primeiro, na borda; o pico térmico ao mesmo tempo ou depois, durante a
descompressão.

\textbf{18.} Sim: o modelo também atinge primeiro seu pico de pressão, e seu pico térmico,
\qty{593}{\celsius} a \qty{7.8}{kbar}, durante a subida, perto das condições da borda
do xisto. O espessamento seguido de uma erosão lenta explica a trajetória do xisto.

\noindent\textbf{Parte IV.}

\textbf{19.} $(623 - 405)/(412.0 - 398.0) = 218/14.0 = \qty{15.6}{\celsius}$ por milhão
de anos.

\textbf{20.} A rocha se resfriou cerca de uma vez e meia mais depressa do que o suposto. Um mineral
que se resfria mais depressa tem menos tempo para perder seu argônio, de modo que se fecha a uma temperatura mais alta:
a verdadeira temperatura de fechamento fica um pouco acima de \qty{405}{\celsius}, e o fechamento
um pouco mais fundo.

\textbf{21.} $(405 - 10)/30 = \qty{13.2}{km}$.

\textbf{22.} $27.5 - 13.2 = \qty{14.3}{km}$ (com as profundidades não arredondadas, 27.52 e 13.17,
\qty{14.35}{km}).

\textbf{23.} $\sqrt{1.5^2 + 1.0^2} = \qty{1.8}{Myr}$: $14.0 \pm \qty{1.8}{Myr}$, logo uma taxa
entre $14.35/15.8 = 0.91$ e $14.35/12.2 = \qty{1.18}{km/Myr}$.

\textbf{24.} $14.35/14.0 = \qty{1.03}{km}$ por milhão de anos, isto é,
\textbf{\qty[text-series-to-math=true]{1.0}{mm.yr^{-1}}}, cerca de trinta vezes mais lenta que
a primeira etapa de \omterm{def:b2:pt-paths:exhumation}{exumação} das rochas alpinas mais profundas, e cerca de duas vezes a taxa de erosão
do modelo de espessamento.
\end{solution}
